\setcounter{footnote}{29}

\printedpage{29}
% head: A. Grothendieck : Dérivateurs, Ch. XVIII, Édition des Universités de Montpellier II et Paris VII
% foot: Version 17 décembre 2025

\section*{4.5\quad Canonical cardinal filtration of a $\mathfrak{U}$-paratopos category}

\textbf{Definition 4.5.1.} By a \textit{$\mathfrak{U}$-paratopos category} we mean a category $\mathcal{M}$ which is a $\mathfrak{U}$-pseudo-topos and accessible, i.e. which satisfies the equivalent conditions of theorem 4.3.6 (p. 70). We say that it is a \textit{$\pi$-$\mathfrak{U}$-paratopos category}, if it is $\pi$-accessible, i.e. if it admits a small generating subcategory whose objects are $\pi$-accessible, or what amounts to the same thing, if % [sic] soient π-accessible -> soient π-accessibles, accord du pluriel
$\mathcal{M}$ is equivalent to a category of the form $\Ind'_{\pi}(\mathcal{M}_{\pi})$, where $\mathcal{M}_{\pi}$ is a small category (or: a category equivalent to a small category), and stable under $\varinjlim$ of cardinality $\leqslant \pi$. (And we can take for $\mathcal{M}_{\pi}$ the full subcategory of $\mathcal{M}$ consisting of the objects
\msp{page 107}
which are $\pi$-accessible).

In what follows, we are going, more generally, to consider a category $C$ equivalent to a small category, stable under $\varinjlim$ of cardinality $\leqslant \pi_0$ (\footnote{There is no need to assume this right away.}), and the category
\[
\mathcal{M} \ \overset{\text{def}}{=}\ \Ind'_{\pi_0}C \ \subset\ C^{\wedge}\ ,
\tag*{(4.5.1)}
\]
where $\pi_0$ is a given infinite cardinal. I propose to introduce on $\mathcal{M}$ a canonical cardinal filtration:

\textbf{Proposition 4.5.2} (\footnote{Is proved (and is true) only if $\pi^{\pi_0} = \pi$, see correction 4.5.6, page 116.}). \textit{Let $F \in \Ob\mathcal{M} \subset \Ob C^{\wedge}$, and let $\pi$ be a cardinal $\geqslant \pi_1 =$} $\sup(\pi_0, \card\Fl C_0)$\textit{, where $C_0$ is a small full subcategory of $C$ such that $C_0 \overset{\sim}{\hookrightarrow} C$. The following conditions are equivalent:}
\begin{enumerate}[leftmargin=3em]
\item[\textit{a)}] $\card\Fl(C_0/F) \leqslant \pi$.
\item[\textit{a$'$)}] $\card\Ob\,(C_0/F) \leqslant \pi$.
\item[\textit{a$''$)}] \textit{For every $x$ in $C_0$,} $\card F(x) \leqslant \pi$.
\item[\textit{b)}] \textit{There exists a small filtering category $I$ (not necessarily ordered),} large relative to $\pi_0$\textit{, with} $\card I \leqslant \pi$\textit{, and a direct system $(x_i)_{i\in I}$ in $C$, such that we have}
\[
F \ \simeq\ \varinjlim_{i\in I} x_i\ ,
\tag*{(4.5.2.1)}
\]
\textit{where the $\varinjlim$ is taken in $\mathcal{M}$, or what amounts to the same thing (by virtue of 4.4.15, $I$ being large relative to $\pi_0$) in $C^{\wedge}$, i.e. $F$ is isomorphic to the ind-representable functor defined by the ind-object $\underline{x} = (x_i)_{i\in I}$ of $C$.}
\item[\textit{c)}] \textit{$F$ is $\pi$-accessible in $\mathcal{M}$.}
\item[\textit{c$'$)}] \textit{$F$ is $\pi$-accessible in $C^{\wedge}$.}
\end{enumerate}
\msp{page 108}

\textsc{Proof}. We have the $\pm$ tautological implications
\[
\begin{tikzcd}[row sep=small]
 & \mathrm{a}' \arrow[r, Rightarrow] & \mathrm{a}'' & \\
\mathrm{a} \arrow[ur, Rightarrow] \arrow[dr, Rightarrow] & & & \\
 & \mathrm{b} \arrow[r, Rightarrow] & \mathrm{c}' \arrow[r, Rightarrow] & \mathrm{c}\ ,
\end{tikzcd}
\tag*{$(*)$}
\]
\printedpage{30}
% head: A. Grothendieck : Dérivateurs, Ch. XVIII, Édition des Universités de Montpellier II et Paris VII
% foot: Version 17 décembre 2025
where a$'$ $\Longrightarrow$ a$''$ comes from the canonical inclusion $F(x) = \Hom_{\mathcal{M}}(x,F) \hookrightarrow \Ob C/F$, a $\Longrightarrow$ b is obtained by taking $I = C_0/F$ and noting that $I$ is stable under direct limits of cardinality $\leqslant \pi_0$ ($C$ being so), hence is large relative to $\pi_0$ (\footnote{We do not need this stability, in any case $C_0/F$ is large relative to $\pi_0$, by virtue of 4.4.3, p. 80.}). We have b $\Longrightarrow$ c$'$ by virtue of the fact that the strictly full subcategory $C^{\wedge}_{\pi}$ of $C^{\wedge}$ defined by
\[
\left\{
\begin{array}{lcl}
C^{\wedge}_{\pi} & = & \text{strictly full subcategory of } C^{\wedge} \text{ consisting of the}\\
 & & F \in \mathcal{M} \text{ which are } \pi\text{-accessible.}
\end{array}
\right.
\tag*{(4.5.2.2)}
\]
is stable under $\varinjlim$ of cardinality $\leqslant \pi$ (SGA 4 I 9.9). Finally c$'$ $\Longrightarrow$ c comes from the fact that the inclusion functor $\mathcal{M} \hookrightarrow C^{\wedge}$ is fully faithful and $\pi_0$-accessible, i.e. commutes with direct limits large relative to $\pi_0$, and a fortiori with those large relative to $\pi$. By the implications of the diagram above, it therefore suffices to prove that we have
% marginal mark (handwritten squiggle) in the left margin of the edition beside the following display
\[
\mathrm{c} \overset{(\mbox{\footnotemark[33]})}{\Longrightarrow} a \qquad\quad \text{and} \qquad\quad \mathrm{a}'' \ \Longrightarrow\ \mathrm{a}\ .
\]
\footnotetext[33]{is proved (and is true) only if $\pi^{\pi_0} = \pi$.}
\setcounter{footnote}{33}

\msp{page 109}

\textit{Let us prove} a$''$ $\Longrightarrow$ a. Let us consider the map
\[
\Fl C/F \ \longrightarrow\ \Fl C
\tag*{$(*)$}
\]
induced by the étale functor $C' = C/F \longrightarrow C$, then the set of arrows above an % [sic] induit -> induite, accord avec « l'application »
arrow $\alpha : x \longrightarrow y$ of $C$ is in one-to-one correspondence with $\Ob C'_y = F(y)$, since such an arrow $\alpha'$ is known when one knows its target, which is an arbitrary object of $F(y)$ ($C' \longrightarrow C$ being étale, i.e. fibrant with discrete fibers). Thus the fiber of $(*)$ at $\alpha$ is of % [sic] le fibre -> la fibre, genre
cardinality $= \card F(y) \leqslant \pi$, so we have
\[
\card\Fl C/F \ \leqslant\ \pi \underbrace{\card\Fl C}_{\leqslant \pi_1} \ \leqslant\ \pi \underbrace{\ \pi_1\ }_{\leqslant \pi} \ \leqslant\ \pi^2 = \pi\ ,
\]
which proves a.

Let us finally prove c $\Longrightarrow$ a. For this, we are going to use the first inclusion $\mathcal{M}_{\pi} \subset \Filt^{\pi}\mathcal{M}$ of SGA 4 I 9.16, which applies as soon as we prove the

% marginal mark (handwritten squiggle) in the left margin of the edition beside Corollaire 4.5.3
\textbf{Corollary 4.5.3} (\footnote{Not proved, condition c) of cardinal filtrations perhaps not satisfied. Probably \textit{false} in general.}). \textit{Under the conditions of 4.5.2, let, for every cardinal $\pi \geqslant \pi_1$,} $\Filt^{\pi}(\mathcal{M})$ \textit{be the strictly full subcategory of $\mathcal{M}$ consisting of the $F \in \Ob\mathcal{M}$ which satisfy condition a) of 4.5.2. Then the family of} $(\Filt^{\pi})_{\pi\geqslant\pi_1}$ \textit{is a}
\msp{page 110}
\textit{cardinal filtration of $\mathcal{M}$.}

Condition a) of loc. cit. 9.12 (the $\uFilt^{\pi}$ are equivalent to small categories) is obvious. Let us prove b), i.e. the stability of $\uFilt^{\pi}$ under filtering $\varinjlim$ $F = \varinjlim F_i$ \textit{large relative to} $\pi_0$, of cardinality $\leqslant \pi$. Since $\mathcal{M} \hookrightarrow C^{\wedge}$ commutes with these limits, for $x$ in $C$
\[
F(x) \ =\ \varinjlim_{i} F_i(x)\ ,
\]
\printedpage{31}
% head: A. Grothendieck : Dérivateurs, Ch. XVIII, Édition des Universités de Montpellier II et Paris VII
% foot: Version 17 décembre 2025
whence $\card F(x) \leqslant \card\Ob I \cdot \pi \leqslant \pi \cdot \pi = \pi$, which (by the equivalence a $\Longleftrightarrow$ a$''$) implies $F \in \Ob\Filt^{\pi}$. Let us prove condition c) of loc. cit. For this, we will need the fact, more precise than b), which we have just established:

\textbf{Corollary 4.5.4} (\footnote{Proved only if $\pi^{\pi_0} = \pi$\,! Probably \textit{false} in general.}). \textit{For every $\pi \geqslant \pi_1$,} $\Filt^{\pi}$ \textit{is stable under $\varinjlim$ (in $\mathcal{M}$) of cardinality $\leqslant \pi$ (not necessarily filtering, a fortiori not necessarily large relative to $\pi_0$).}

It suffices, by virtue of what is already proved of of lemma 4.4.14 (\footnote{We need a variant of lemma 4.4.14, for the $\varinjlim$ of cardinality $\leqslant \pi'$, which is proved in the same way.}), to show that $\Filt^{\pi}$ % [sic] de du lemme -> du lemme, mot répété
is stable under $\varinjlim$ \textit{of cardinality} $\leqslant \pi_0$. If $\varphi : J \longrightarrow \Filt^{\pi}$ is a functor with $\card J \leqslant \pi_0$, we put it in indexed form
\[
\psi\ :\ I \times J \ \longrightarrow\ C\ ,
\]
thanks to 4.14.3.b, with $I$ large relative to $\pi_0$.

\msp{page 111}

I recall that we could take
\[
I \ =\ \uHom(J,C)/\varphi\ ,
\]
where $\uHom(J,C) \hookrightarrow \uHom(J, \mathcal{M} = \Ind'_{\pi}C)$, at the cost of not requiring $I$ to be ordered. We would like to prove that $\card\Fl I \leqslant \pi$. The trouble is that we only arrive at the upper bound
\[
\card\Fl I \ \leqslant\ \pi^{\pi_0}\ .
\]
(Even if $J$ is discrete, we do no better, since then
\[
I \ \simeq\ \prod_{j\in J} C/\varphi(j)\ , \qquad \text{whence} \quad \Fl I = \prod_{j\in J} \Fl C/\varphi(j)\ ,
\]
and if $\card\Fl C/\varphi(j) = \pi$ for every $j$, we find $\card\Fl I = \pi^{\pi_0}$ -- there is no way to do better\,!)

So I cannot manage to prove 4.5.4 directly, which would have come in handy for proving condition c) of loc. cit. But to tell the truth, we have done here the work of the proof of theorem 4.3.6, part b $\Longrightarrow$ c -- we have already proved there that the filtration we have just constructed (by condition a) of 4.5.2) is a cardinal filtration. Condition c) of SGA 4 I 9.12 is proved on page 76. But I am going to redo
\msp{page 112}
the proof more seriously.

So let $F \in \Ob\mathcal{M}$, consider
\[
J \ =\ C/F\ ,
\]
it is a category stable under $\varinjlim$ of cardinality $\leqslant \pi_0$, hence large relative to $\pi_0$, and we have
\[
F \ \simeq\ \varinjlim_{j\in J} x_j \qquad\quad (\,\varinjlim \ \text{in } \mathcal{M})
\]
\printedpage{32}
% head: A. Grothendieck : Dérivateurs, Ch. XVIII, Édition des Universités de Montpellier II et Paris VII
% foot: Version 17 décembre 2025
denoting by $j \longmapsto x_j$ the canonical functor $C/F \longrightarrow C$. Then let $I$ be the set of subcategories $J'$ of $J$ which are \textit{of cardinality $\leqslant \pi$ and large relative to $\pi_0$, $I$} ordered by inclusion, and for every $J' \in I$, let us put
\[
X_{J'} \ =\ \varinjlim_{j\in J'} x_j\ .
\]
As $J'$ is large relative to $\pi_0$, the direct limit in $\mathcal{M}$ is also a limit in $C^{\wedge}$, whence it follows at once, since $\card\Ob J' \leqslant \card\Fl J' \leqslant \pi$, that $X_{J'}$ is in $\Filt^{\pi}$. Then
\[
J' \ \longmapsto\ X_{J'} \ \in\ \Ob\Filt^{\pi}
\]
is a functor
\[
I \ \longrightarrow\ \Filt^{\pi}\ ,
\]
whose direct limit is that of the $x_j$, hence isomorphic to $F$. On the other hand, it is immediate that $I$ is large relative to $\pi$, i.e. that if $(J'_{\alpha})_{\alpha\in A}$ is
\msp{page 113}
a family of elements of $I$, it is bounded above by a $J' \in I$. This follows from the

% marginal mark (handwritten squiggle) in the left margin of the edition beside Lemme 4.5.5
\textbf{Lemma 4.5.5} (\footnote{Is proved \textit{only} if $\pi^{\pi_0} = \pi$\,! Probably false in general.}). \textit{Let $\pi$, $\pi_0$, with $\pi > \pi_0$, be infinite cardinals, and let $J$ be a category large relative to $\pi_0$. Then for every subset $\Phi$ of} $\Fl J$\textit{, such that} $\card\Phi \leqslant \pi$\textit{, there exists a subcategory $J'$ of $J$} large relative to $\pi_0$\textit{, such that} $\card J' \leqslant \pi$ (\footnote{This is OK if we require only $\card J' \leqslant \pi^{\pi_0}$.})\textit{, and} $\Fl J' \supset \Phi$.

This lemma being admitted, we have therefore represented an arbitrary object $F$ of $\mathcal{M}$ as a $\varinjlim\limits_{i\in I} X_i$, with $I$ an ordered set large relative to $\pi$, the $X_i$ in $\Filt^{\pi}$. It remains to prove that if $\pi' > \pi$ and $F \in \Filt^{\pi'}$, then we can take $I$ such that $\card I \leqslant \pi'^{\pi}$. But $F \in \Filt^{\pi'}$ means that $J = C/F$ satisfies $\card\Fl J \leqslant \pi'$. Then the set of subsets of cardinality $\leqslant \pi$ of $\Fl J$ has a cardinality bounded above by $\pi'^{\pi}$, a fortiori $\card I \leqslant \pi'^{\pi}$, which finishes proving that $(\Filt^{\pi})_{\pi\geqslant\pi_1}$ is a cardinal filtration on $\mathcal{M}$.

\msp{page 114}

To prove the lemma, we begin by saturating $\Phi$ under finite composition, we find a subcategory $J'_0$ of $J$ such that $\card\Fl J'_0 \leqslant \pi_0$ \ldots{} But I notice that to prove the lemma, I again need $\pi^{\pi_0} = \pi$ -- an envelope of $J'$ large relative to $\pi_0$ is going to have a cardinality bounded above (at best) by $\pi^{\pi_0}$.

Let us take e.g.
\[
J \ =\ \mathfrak{P}(S) \qquad\quad \text{(ordered set)}
\]
with $\card S = \pi > \pi_0$, so $J$ is large relative to $\pi_0$. Let us consider the map
\[
S \ \longrightarrow\ J = \mathfrak{P}(S)\ , \qquad\quad s \longmapsto \{s\}\ ,
\]
and let us take for $\Phi \subset \Fl J$ the set of $\mathrm{id}_{\{s\}}$. Thus we are looking for a subset $J'$ of $J$, containing $S$, and large relative to $\pi_0$ for the order induced by $J$. We see at once that the smallest subset $J'$ that does the job is the set $\mathfrak{P}_{\pi_0}(S)$ of subsets of $S$ of cardinality $\leqslant \pi_0$. But we find
\[
\card J' \cdot \pi_0! \ =\ \pi^{\pi_0}\ ,
\]
\printedpage{33}
% head: A. Grothendieck : Dérivateurs, Ch. XVIII, Édition des Universités de Montpellier II et Paris VII
% foot: Version 17 décembre 2025
\msp{page 115}
where
\[
\begin{aligned}
\pi_0! \ &=\ \card\Aut\,(E_0) \qquad\quad (E_0 \text{ a set of cardinality } \pi_0)\\
&=\ \pi_0^{\pi_0}\\
&=\ 2^{\pi_0}\ ,
\end{aligned}
\]
so
\[
\card J' \cdot 2^{\pi_0} \ =\ \pi^{\pi_0}\ .
\tag*{$(*)$}
\]
On the other hand, it is immediate that $\card J' \geqslant \card\mathfrak{P}(E_0) = 2^{\pi_0}$ (where $E_0$ is a subset of $S$ of cardinality $\pi_0$), so the left-hand side of $(*)$ is equal to $\card J'$, so we find
\[
\card J' \ =\ \pi^{\pi_0}\ ,
\]
sorry\,!

Finally, I have convinced myself that proposition 4.5.2, and the preceding corollaries and lemmas by which I tried to establish it, is true only for those cardinals $\pi > \pi_0$ such that
\[
\pi^{\pi_0} \ =\ \pi\ ,
\]
which, assuming the generalized continuum hypothesis, is equivalent to
\[
\pi^{\pi_0} \ <\ 2^{\pi}\ (\ =\ \pi^{\pi}\ )
\]
(or again to $\pi^{\pi_0} \neq 2^{\pi}$ $(= \pi^{\pi})$). I am going to give
\msp{page 116}
here a corrected statement.

\textbf{Proposition 4.5.6.}
\begin{enumerate}[leftmargin=3em]
\item[\circled{1\textsuperscript{o}}] \textit{Under the general hypotheses of 4.5.2 on $C$, $\pi_0$, $\pi_1$, $\pi$, the conditions a, a$'$, a$''$, b, c$'$ of this proposition are equivalent (provided $\pi \geqslant \pi_1$)} (\footnote{\textbf{N.B.} For this equivalence, the hypothesis that $C$ be stable under $\varinjlim$ of cardinality $\leqslant \pi$ is unnecessary, and also for c$'$) to imply c), and to be equivalent to it if $\pi = \pi^{\pi_0} \geqslant \pi_1$, cf. 4.7.1, p. 161.})\textit{, and imply condition c ($F$ is $\pi$-accessible). When $\pi^{\pi_0} = \pi$ (in particular when $\pi = 2^{c}$, with $c \geqslant \pi_0$), then all the conditions a to c$'$ of 4.5.2 are equivalent.}
\item[\circled{2\textsuperscript{o}}] \textit{More generally and more precisely, for $\pi \geqslant \pi_1$ given, let us consider the conditions on $\pi$:}
\begin{enumerate}[leftmargin=2em]
\item[\textit{1)}] $\pi^{\pi_0} = \pi$.
\item[\textit{2)}] $\Filt^{\pi}(\mathcal{M})$ \textit{(defined in 4.5.3) is stable under $\varinjlim$ of cardinality $\leqslant \pi$.}
\item[\textit{3)}] \textit{$\forall$ $F$ in $\mathcal{M}$, we have $F \simeq \varinjlim\limits_{I} X_i$, with $I$ an ordered set large relative to $\pi$, the $X_i$ in} $\Filt^{\pi}(\mathcal{M})$\textit{. Moreover, if $\pi' > \pi$ and $F \in$} $\Filt^{\pi'}(\mathcal{M})$\textit{, we can take $I$ of cardinality $\leqslant \pi'^{\pi}$.}
\item[\textit{4)}] \textit{As in 3), without the bound on cardinalities, and assuming only $I$ a category large relative to $\pi$ (not necessarily an ordered category).} % [sic] catégorie ordonné -> catégorie ordonnée, accord
\item[\textit{5)}] \textit{Condition c) of 4.5.2 implies the other conditions (they are equivalent, by 1°), in other words}
\[
\mathcal{M}_{\pi} \ \subset\ \Filt^{\pi} \qquad\quad \textit{(so } \mathcal{M}_{\pi} = \Filt^{\pi}\textit{, since we know that } \Filt^{\pi} \subset \mathcal{M}_{\pi}\textit{)}\ ,
\]
\end{enumerate}
\end{enumerate}